\documentclass[11pt,a4paper]{article}
\usepackage[margin=2.2cm]{geometry}
\usepackage{fontspec}
\setmainfont{TeX Gyre Pagella}
\usepackage{enumitem}
\usepackage{xcolor,listings}
\lstset{basicstyle=\ttfamily\small,frame=single,rulecolor=\color{gray!50},columns=fullflexible,keepspaces=true,breaklines=true}
\usepackage[hidelinks]{hyperref}
\setlist{nosep}
\setlength{\parskip}{4pt}

\title{Competitive Programming and Contests\\Hands-On 2 -- Report}
\author{Gabriele Marsili}
\date{November 2026}

\begin{document}\sloppy
\maketitle

Both solutions live in \texttt{lib.rs}, which contains \texttt{solve\_p1}, \texttt{solve\_p2}, the two segment
trees they rely on, the parsing of the input and the tests, while \texttt{main.rs} does nothing more than read
the input and print the answers, one per line. Since \texttt{main.rs} imports everything from the library, the
package has to be called \texttt{handson2}:
\begin{lstlisting}
cargo run --release -- 1 < testsets/p1/input0.txt     # Problem 1
cargo run --release -- 2 testsets/p2/input0.txt       # Problem 2, input from a file
\end{lstlisting}
The input can come from stdin, from a file whose path is passed as an argument, or directly from the command
line. Throughout the report, $n$ is the size of the array (or the number of segments) and $m$ the number of
queries.

\section*{Problem 1 -- Min and Max}

\textbf{Indices.} The queries in the statement and in the test set are 1-based, so \texttt{solve\_p1}
(line 213) shifts them to 0-based before passing them to the tree.

\textbf{Data structure.} I used a segment tree (\texttt{SegmentTree}, line 100) stored in an array, with the
root at index 1 and the children of node $k$ at $2k$ and $2k+1$; every node covers a range $A[lo..hi]$ and
keeps the maximum of that range. An array of $4n$ entries is always large enough, and the tree is built
recursively, computing each node from its two children, in $O(n)$.

\textbf{Max$(i,j)$.} The query is the usual recursive visit: a node that lies outside $[i,j]$ returns
$-\infty$ (\texttt{i64::MIN}), a node fully contained in $[i,j]$ returns its value without going further, and
only a node that partially overlaps the range asks both children and keeps the larger of the two answers.

The reason why this costs $O(\log n)$ is that, on each level of the tree, at most two nodes can partially
overlap $[i,j]$. A node partially overlaps the range only if one of the two endpoints falls strictly inside
it, i.e.\ it contains both $i-1$ and $i$, or both $j$ and $j+1$, and since the nodes of a level are disjoint
there is at most one node containing $i$ and at most one containing $j$. Nodes outside or fully inside the
range stop the recursion, so the nodes visited on the next level are only the children of these (at most)
two partial nodes, which means at most four nodes per level and $O(\log n)$ nodes over the $\lceil\log_2
n\rceil + 1$ levels. The same argument holds for the update described below, which stops in the same cases.

\textbf{Update$(i,j,T)$.} Updating each position on its own would cost $O((j-i+1)\log n)$, that is
$O(n \log n)$ per query in the worst case, so to bring it down to $O(\log n)$ I used lazy propagation. What
makes it work here is that taking the minimum with $T$ commutes with the maximum:
\[ \max_{p \in [lo,hi]} \min(A[p], T) = \min\Big(\max_{p \in [lo,hi]} A[p],\ T\Big). \]
This means that when the range of a node is fully inside $[i,j]$ its new value is simply $\min(t[k], T)$
and the recursion can stop right there, while the children, which are left untouched, are informed later
through \texttt{lazy[k]}, where the pending value $T$ is stored (\texttt{apply}, line 142). One value per
node is enough, because two pending updates on the same node merge into one, as
$\min(\min(x, T_1), T_2) = \min(x, \min(T_1, T_2))$.

Whenever a query or an update needs to go below a node, \texttt{push} (line 152) first hands the pending value
down to the two children and clears it, which is also why \texttt{max} takes \texttt{\&mut self}. A leaf is
always either fully inside or fully outside a range, so \texttt{push} is never called on a leaf. Once the
recursive calls return, the update recomputes the node as the maximum of its children (line 182).

\textbf{Observation:} if $t[k] \le T$, every value in the range of the node is already at most $T$, so
nothing would change and the update stops there without recording anything.

\textbf{Complexity:}
\begin{itemize}
\item $O(n)$ to build the tree and $O(\log n)$ for each update and each query, so $O((n+m)\log n)$ in total;
\item $O(n)$ space, for the $4n$ values and the $4n$ pending updates.
\end{itemize}

\section*{Problem 2 -- Is There}

Here positions and queries are 0-based. If $\mathit{cover}[p]$ is the number of segments that contain
position $p$, the query IsThere$(i,j,k)$ asks whether $\mathit{cover}[p] = k$ for at least one $p \in [i,j]$.

\textbf{Step 1: coverage.} Since the segments never change, $\mathit{cover}$ can be computed once with a
difference array (line 303): for each segment $\langle l, r\rangle$ I add $1$ at $l$ and $-1$ at $r+1$, and
the prefix sums give $\mathit{cover}[p]$ in $O(n)$. In the same pass I group the positions by their
coverage, so that \texttt{pos\_with[k]} is the list of the positions $p$ with $\mathit{cover}[p] = k$
(line 312); with $n$ segments, $0 \le \mathit{cover}[p] \le n$.

\textbf{Step 2: queries offline, grouped by $k$.} Neither a maximum nor a sum over $\mathit{cover}$ can tell
whether the value $k$ appears in a range, so I answer the queries offline, grouping them by $k$ and using a
segment tree over the positions (\texttt{CountTree}, line 230) in which every node stores how many positions of
its range are \emph{active}, with point updates and range sums. Then, for each $k = 0, \dots, n$ (line 324):
\begin{enumerate}
\item I activate the positions in \texttt{pos\_with[k]};
\item each query with this $k$ is answered 1 if and only if $[i,j]$ contains at least one active position;
\item I deactivate the same positions, so that the tree is empty again for the next value of $k$.
\end{enumerate}
The values of $k$ that have no queries are skipped, and a query with $k > n$ is answered 0 straight away,
since no position can be covered by more than $n$ segments. The answers are stored by query index, so they
come out in the same order as the input.

\textbf{Complexity:}
\begin{itemize}
\item every position is activated and deactivated at most once ($O(n\log n)$), every query is a single range
sum ($O(m \log n)$) and the grouping is $O(n+m)$, so $O((n+m)\log n)$ in total;
\item $O(n+m)$ space, for the coverage, the two groupings and the tree with $4n$ nodes.
\end{itemize}

\section*{Tests}

All the tests pass, the code is formatted with \texttt{cargo fmt} and \texttt{cargo clippy} reports no
warnings.
\begin{itemize}
\item Official test set: a separate test runner, not included in the submission, compares the output of the
binary with every \texttt{outputK.txt}, and all 11 files of Problem 1 and all 8 of Problem 2 match.
\item Unit tests in \texttt{lib.rs}: the examples of the statement for both problems, an update covering the
whole array followed by a query on a single position (to check that the pending value stored at the root
actually reaches the leaf), and queries with $k > n$.
\item Randomized tests: each solution is compared with a brute-force version on 3\,000 random inputs, with a
small $n$ so that every branch of the recursion gets exercised.
\item Parsing: the examples of the statement, the whole input on a single line, truncated input and
non-numeric tokens.
\item Performance: on random inputs with $n = m = 200\,000$, where for Problem 1 half of the queries are
updates on the whole array, the release build takes about 0.2 s for Problem 1 and 0.2--0.3 s for Problem 2,
input reading included.
\end{itemize}

\section*{Sources}
\begin{itemize}
\item The statement of the Hands-On: \url{https://pages.di.unipi.it/rossano/blog/2024/handson2/}.
\item The lectures of the course on segment trees and lazy propagation.
\item CP-Algorithms, ``Segment Tree'', in particular the sections on range updates and lazy propagation:
\url{https://cp-algorithms.com/data_structures/segment_tree.html}.
\end{itemize}

\section*{Use of AI}
I designed the structure of the two solutions (segment tree with lazy propagation for Problem 1, coverage
with a difference array and offline queries grouped by $k$ for Problem 2), with Claude Code reviewing it.
The code was written by Claude Code; I reviewed it carefully and reworked some parts. Claude Code also
generated \texttt{main.rs}, the input parsing and the tests, and fixed the remaining bugs. Claude Code also
wrote this report, following my indications and corrections, and I reviewed the final version.

\end{document}
